% ============================================================
% Section 2: Model
% Journal of Financial Econometrics style
% ============================================================

\section{Model}\label{sec:model}

\subsection{Price process and observation scheme}

Let $(X(t),Y(t))$ be the bivariate latent log-price process on $[0,1]$,
assumed to satisfy
\begin{align}\label{eq:bivariateBM}
  dX(t) &= \sigma_1\, dB_1(t), \nonumber \\
  dY(t) &= \sigma_2\, dB_2(t),
\end{align}
where $B_1,B_2$ are standard Brownian motions with instantaneous correlation
$\rho$, so $\langle dB_1, dB_2 \rangle = \rho\,dt$, and $\sigma_1,\sigma_2 > 0$,
$\rho \in (-1,1)$ are constants.  The target parameter is the integrated
covariance
\begin{equation}\label{eq:target}
  \sigma_{12} \;=\; \sigma_1 \sigma_2 \rho.
\end{equation}

Asset $X$ is observed at times $\mathcal{T} = (t_0, t_1, \ldots, t_m)$ and
asset $Y$ at times $\mathcal{U} = (u_0, u_1, \ldots, u_n)$, with
$t_0 = u_0 = 0$ and $t_m = u_n = 1$.  Observation times are fixed and
independent of the price process; the two grids need not share any points in
common.  The returns are
\begin{equation*}
  R_i = X(t_i) - X(t_{i-1}), \quad i = 1, \ldots, m, \qquad
  S_j = Y(u_j) - Y(u_{j-1}), \quad j = 1, \ldots, n.
\end{equation*}
We write $\tau_i^X = t_i - t_{i-1}$ and $\tau_j^Y = u_j - u_{j-1}$ for the
lengths of the return intervals.

\subsection{Overlap structure}

The \emph{overlap length} between the $i$th return interval of $X$ and the
$j$th return interval of $Y$ is
\begin{equation}\label{eq:overlap}
  \tau_{ij} \;=\; \bigl|(t_{i-1},\, t_i] \cap (u_{j-1},\, u_j]\bigr|,
\end{equation}
where $|\cdot|$ denotes Lebesgue measure.  We have $\tau_{ij} \geq 0$, with
$\tau_{ij} = 0$ whenever the two intervals are disjoint.

The key consequences of the model~\eqref{eq:bivariateBM} are
\begin{align}\label{eq:moments}
  \mathbb{E}[R_i]     &= 0, &\quad \mathbb{E}[R_i^2] &= \sigma_1^2\,\tau_i^X, \nonumber \\
  \mathbb{E}[S_j]     &= 0, &\quad \mathbb{E}[S_j^2] &= \sigma_2^2\,\tau_j^Y, \\
  \mathbb{E}[R_i S_j] &= \sigma_{12}\,\tau_{ij}. & & \nonumber
\end{align}
The cross-moment identity in~\eqref{eq:moments} follows from the It\^{o}
isometry: $\mathbb{E}[R_i S_j] = \rho\,\sigma_1\sigma_2\,|J_i^X \cap J_j^Y|
= \sigma_{12}\,\tau_{ij}$.  It is the sole probabilistic input to all
subsequent results; everything downstream is algebra.

Further, since increments of Brownian motion over non-overlapping intervals
are independent and jointly Gaussian, the Isserlis--Wick theorem gives
\begin{equation}\label{eq:isserlis}
  \mathrm{Cov}(R_i S_j,\, R_k S_l)
  \;=\;
  \sigma_1^2\sigma_2^2\,\tau_i^X\tau_j^Y\,\mathbf{1}_{(i,j)=(k,l)}
  \;+\;
  \sigma_{12}^2\,\tau_{il}\,\tau_{kj},
\end{equation}
for all $(i,j),(k,l) \in \{1,\ldots,m\}\times\{1,\ldots,n\}$.
Equation~\eqref{eq:isserlis} is the foundation for the exact
finite-sample variance formula in Theorem~\ref{thm:genHY}.

\begin{remark}[Constant parameters]
The assumption of constant $\sigma_1,\sigma_2,\rho$ is made for notational
clarity.  All results hold for the time-varying case with
$\sigma_1(t),\sigma_2(t),\rho(t)$ by replacing $\sigma_{12}$ with
$\int_0^1 \sigma_1(t)\sigma_2(t)\rho(t)\,dt$ throughout, since the
cross-moment identity $\mathbb{E}[R_i S_j] = \int_{J_i^X \cap J_j^Y}
\sigma_{12}(t)\,dt$ continues to hold under It\^{o}'s isometry.
\end{remark}

\begin{remark}[No microstructure noise]
We assume noise-free observations of the efficient log-price.  Extensions
to the additive noise model, following \citet{zhang2011estimating} and
\citet{bibinger2014estimating}, are discussed in
Section~\ref{sec:discussion}.
\end{remark}